\documentclass{article}
\title{Practice Problem 2.28}
\author{CSAPP}
\usepackage{booktabs, amssymb, parskip, amsmath, tabularx}

\begin{document}
    \maketitle
    \begin{flushleft}
        We can represent a bit pattern of of length $w = 4$
        with a single $hex$ digit. For unsigned representation of these
        digit, using equation: \\

        \[
         -^u_w x = \begin{cases}
         x,       & x = 0\\
         2^w - x  & x > 0 
        \end{cases}
        \]
        to fill in the following table giving the values and the bit representations
        $(in hex)$ of the unsigned additive inverses of the digits shown
        \\

        \begin{tabular}{cc|cc}
        \toprule
        \multicolumn{2}{c|}{$x$} & \multicolumn{2}{c}{$4 -^u x$} \\
        \midrule
        Hex & Decimal & Decimal & Hex \\
        \midrule
        $1$ & $1$  & $15$ & $F$\\
        $4$ & $4$  & $12$ & $C$\\
        $7$ & $7$  & $9$  & $9$\\
        $A$ & $10$ & $6$  & $6$\\
        $E$ & $14$ & $2$  & $2$\\
        \bottomrule
        \end{tabular}
    \end{flushleft}
\end{document}